\documentclass[11pt, a4paper]{article}

\usepackage[utf-8]{inputenc}
\usepackage[T1]{fontenc}
\usepackage{lmodern}
\usepackage{amsmath, amssymb, amsfonts}
\usepackage{geometry}
\usepackage{graphicx}
\usepackage{hyperref}
\usepackage{booktabs}
\usepackage{color}
\usepackage{cite}
\usepackage{microtype}
\usepackage{xcolor}
\usepackage{array}
\usepackage{longtable}

\geometry{
    a4paper,
    left=25mm,
    right=25mm,
    top=25mm,
    bottom=25mm
}

\hypersetup{
    colorlinks=true,
    linkcolor=blue,
    filecolor=magenta,
    urlcolor=cyan,
    pdftitle={KPP K-Profile Parameterization: Comprehensive Technical Report},
    pdfauthor={Ocean Physics Research Division}
}

\title{\textbf{KPP K-Profile Parameterization} \\ \large Comprehensive Technical Report}
\author{Ocean Physics Research Division}
\date{\today}

\begin{document}

\maketitle

\begin{abstract}
This report provides a comprehensive technical analysis of the KPP (K-Profile Parameterization) vertical mixing scheme as implemented in MITgcm and used globally in ocean circulation models. KPP is a diagnostic boundary-layer mixing scheme based on Large et al.\ (1994) that diagnoses a boundary-layer depth from bulk Richardson number criterion and then applies enhanced mixing coefficients within the boundary layer while maintaining background interior mixing. The report documents the scientific foundation, package structure, complete computational sequence, parameter definitions, inputs and outputs, integration with the ocean model, and comparison with the alternative GGL90 (TKE closure) scheme. This documentation serves both as a reference for model users and as a guide for implementation in research and operational systems.
\end{abstract}

\tableofcontents
\newpage

\section{Purpose and Scope}

\subsection{What is KPP?}

KPP (K-Profile Parameterization) is a diagnostic boundary-layer mixing scheme developed by Large et al.\ (1994). It computes vertical mixing coefficients at each ocean column and model time step by:

\begin{enumerate}
    \item Diagnosing the depth of the oceanic boundary layer from a bulk Richardson number criterion
    \item Computing enhanced mixing within the diagnosed boundary layer using shape functions
    \item Maintaining background interior mixing outside the boundary layer
    \item Optionally computing nonlocal transport coefficients for tracers in unstable conditions
\end{enumerate}

KPP is called ``diagnostic'' because it does not solve a prognostic equation; instead, it computes mixing coefficients directly from the current model state.

The scheme outputs three primary diagnostic fields per column:

\begin{itemize}
    \item \texttt{KPPviscAz}: Vertical eddy viscosity coefficient for momentum $K_M(z)$ [m$^2$/s]
    \item \texttt{KPPdiffKzT}: Vertical diffusion coefficient for temperature $K_T(z)$ [m$^2$/s]
    \item \texttt{KPPdiffKzS}: Vertical diffusion coefficient for salinity/tracers $K_S(z)$ [m$^2$/s]
    \item \texttt{KPPghat}: Nonlocal transport coefficient $\gamma(z)$ [s/m$^2$]
    \item \texttt{KPPhbl}: Diagnosed boundary-layer depth $h_{bl}$ [m]
\end{itemize}

\subsection{Key Physics Features}

\begin{description}
    \item[Boundary Layer Diagnosis] Identifies surface-forced mixed layer from balance of shear and buoyancy
    \item[Shape Functions] Uses cubic/polynomial profiles to represent enhanced BL mixing
    \item[Interior Mixing] Background mixing from stratification and shear instability
    \item[Nonlocal Transport] Explicit $\gamma(z)$ term for convective fluxes
    \item[Surface Forcing Coupling] Directly couples wind, heat flux, and freshwater to boundary layer depth
    \item[Rapid Adjustment] Instantaneous response (no memory); suitable for initialization
\end{description}

\subsection{Applications}

KPP is used in:

\begin{itemize}
    \item Subtropical and tropical ocean simulations (designed for shallow boundary layers)
    \item Climate models requiring fast initialization
    \item Global state estimation where rapid data assimilation is needed
    \item Coupled atmosphere-ocean systems
    \item Research applications requiring transparent boundary layer physics
\end{itemize}

\section{Scientific Background and Core Equations}

\subsection{Boundary Layer Depth Diagnosis}

The fundamental quantity in KPP is the mixed-layer depth $h_{bl}$, diagnosed from the bulk Richardson number:

\begin{equation}
\text{Ri}_b(z) = \frac{\Delta b(z) \cdot (z_{ref} - z)}{|\Delta \mathbf{V}(z)|^2 + V_t^2(z)}
\end{equation}

where:

\begin{itemize}
    \item $\Delta b(z) = g[\rho(z) - \rho_{ref}]/\rho_0$ — Buoyancy difference from surface [m/s$^2$]
    \item $\Delta \mathbf{V}(z) = |\mathbf{u}(z) - \mathbf{u}_{ref}|$ — Horizontal velocity difference from reference [m/s]
    \item $V_t(z)$ — Turbulent velocity scale [m/s]
    \item $z_{ref}$ — Reference depth (typically surface)
\end{itemize}

The boundary layer depth $h_{bl}$ is defined as the shallowest depth where:

\begin{equation}
\text{Ri}_b(h_{bl}) = \text{Ri}_{cr} = 0.3
\end{equation}

This criterion captures the balance between destabilizing shear production and stabilizing buoyancy stratification. When shear dominates ($\text{Ri}_b$ is small), $h_{bl}$ is deep. When stratification dominates ($\text{Ri}_b$ is large), $h_{bl}$ is shallow.

\subsection{Mixing Coefficients: Interior Region ($z > -h_{bl}$)}

Interior mixing (outside the boundary layer) consists of three components:

\begin{equation}
K(z) = K_{bg} + K_{shear} + K_{conv}
\end{equation}

\subsubsection{Background Diffusivity}

A constant baseline from internal wave activity:

\begin{equation}
K_{bg} = \text{difm0} = 5 \times 10^{-3} \text{ m}^2/\text{s}
\end{equation}

\subsubsection{Shear Instability}

Local Richardson number controls shear-driven mixing:

\begin{equation}
\text{Ri}(z) = \frac{N^2(z)}{S^2(z) + \phi_\epsilon}
\end{equation}

where $N^2 = -g/\rho_0 \cdot \partial\rho/\partial z$ and $S^2 = (\partial u/\partial z)^2 + (\partial v/\partial z)^2$.

Shear mixing coefficient:

\begin{equation}
K_{shear}(z) = K_0 \cdot f(\text{Ri})
\end{equation}

where:

\begin{equation}
f(\text{Ri}) = \left[1 - \min\left(\frac{\text{Ri}}{\text{Ri}_\infty}, 1\right)\right]^3
\end{equation}

with $K_0 = 5 \times 10^{-3}$ m$^2$/s and $\text{Ri}_\infty = 0.7$.

\subsubsection{Static Instability}

When $N^2 < 0$ (density inversion), apply enhanced diffusivity:

\begin{equation}
K_{conv}(z) = \text{difmcon} = 0.1 \text{ m}^2/\text{s} \quad \text{if} \quad N^2 < 0
\end{equation}

This parameterizes gravitational overturning in unstable layers.

\subsection{Mixing Coefficients: Boundary Layer ($z < -h_{bl}$)}

Within the boundary layer, mixing is parameterized via shape functions:

\begin{equation}
K(z) = h_{bl} \cdot w(z) \cdot G(z; h_{bl}) \cdot [1 + \sigma \cdot G_1(\sigma)]
\end{equation}

where:

\begin{itemize}
    \item $\sigma = -z/h_{bl}$ — Normalized depth within boundary layer (0 at surface, 1 at $h_{bl}$)
    \item $w(z)$ — Turbulent velocity scale (differs for momentum vs.\ scalars)
    \item $G(\sigma)$ — Shape function (cubic polynomial)
    \item $G_1$ — Boundary condition matching at $z = -h_{bl}$
\end{itemize}

\subsubsection{Turbulent Velocity Scales}

For unstable conditions ($B_0 < 0$, convection), scales are computed from lookup tables:

\begin{equation}
\zeta = \frac{\kappa h_{bl} |B_0|}{u_*^3}
\end{equation}

where $u_* = (|\boldsymbol{\tau}|/\rho_0)^{1/4}$ is friction velocity and $B_0$ is surface buoyancy forcing.

For stable conditions ($B_0 > 0$), direct formula applies.

\subsubsection{Shape Functions}

The cubic shape function ensures:

\begin{itemize}
    \item Smooth transition from surface (enhanced mixing) to boundary layer base
    \item Continuous matching with interior mixing at $z = -h_{bl}$
    \item Proper boundary conditions from turbulent energy balance
\end{itemize}

\subsection{Nonlocal Transport Coefficient}

For tracer (temperature, salinity) fluxes, KPP includes a nonlocal transport term:

\begin{equation}
\left\langle w'c' \right\rangle = -K_H \frac{\partial c}{\partial z} + w_e c^* \gamma(z)
\end{equation}

where:

\begin{itemize}
    \item First term — Standard diffusive flux
    \item $\gamma(z)$ — Nonlocal transport coefficient [s/m$^2$]
    \item $c^*$ — Cumulative surface forcing effect
    \item $w_e$ — Entrainment velocity at boundary layer base
\end{itemize}

The nonlocal coefficient $\gamma(z)$ is:

\begin{equation}
\gamma(z) = -\frac{cstar}{\int_0^{-h_{bl}} G(\sigma) d\sigma} \cdot G(\sigma) \quad \text{for} \quad z < -h_{bl}
\end{equation}

where $cstar$ is a dimensionless constant (typically 10).

Physically, $\gamma$ represents the non-local effect of large coherent eddies transporting properties vertically without local gradients (e.g., plume dynamics in convection).

\section{Package Structure}

\subsection{Directory Organization}

The KPP package is located in \texttt{pkg/kpp/} of the MITgcm repository and contains approximately 3000 lines of Fortran code organized as:

\subsubsection{Core Computation}

\begin{itemize}
    \item \texttt{kpp\_calc.F} — Main entry point and orchestration (560 lines)
    \item \texttt{kpp\_routines.F} — 1D column solver routines (1090 lines)
    \item \texttt{kpp\_forcing\_surf.F} — Surface forcing computation (203 lines)
\end{itemize}

\subsubsection{Supporting Physics}

\begin{itemize}
    \item \texttt{kpp\_calc\_visc.F} — Viscosity coefficient extraction (64 lines)
    \item \texttt{kpp\_calc\_diff\_t.F} — Temperature diffusivity extraction (52 lines)
    \item \texttt{kpp\_calc\_diff\_s.F} — Salinity diffusivity extraction (52 lines)
    \item \texttt{kpp\_calc\_diff\_ptr.F} — Passive tracer diffusivity extraction (64 lines)
    \item \texttt{kpp\_transport\_t.F} — Temperature nonlocal transport (94 lines)
    \item \texttt{kpp\_transport\_s.F} — Salinity nonlocal transport (94 lines)
    \item \texttt{kpp\_transport\_ptr.F} — Passive tracer nonlocal transport (94 lines)
\end{itemize}

\subsubsection{Optional Physics}

\begin{itemize}
    \item \texttt{kpp\_doublediff.F} — Salt fingering and diffusive convection (178 lines)
\end{itemize}

\subsubsection{Initialization and Configuration}

\begin{itemize}
    \item \texttt{kpp\_readparms.F} — Read data.kpp namelist (227 lines)
    \item \texttt{kpp\_init\_fixed.F} — Fixed initialization (221 lines)
    \item \texttt{kpp\_init\_varia.F} — Variable initialization (108 lines)
    \item \texttt{kpp\_check.F} — Configuration validation (116 lines)
\end{itemize}

\subsubsection{I/O and Diagnostics}

\begin{itemize}
    \item \texttt{kpp\_output.F} — Snapshot output (98 lines)
    \item \texttt{kpp\_diagnostics\_init.F} — Diagnostic field registration (198 lines)
\end{itemize}

\subsubsection{Parallel Support}

\begin{itemize}
    \item \texttt{kpp\_do\_exch.F} — MPI halo exchanges (68 lines)
\end{itemize}

\subsubsection{Configuration Headers}

\begin{itemize}
    \item \texttt{KPP.h} — Package state declarations (178 lines)
    \item \texttt{KPP\_PARAMS.h} — Parameter definitions (241 lines)
    \item \texttt{KPP\_OPTIONS.h} — Compile-time CPP options (85 lines)
\end{itemize}

\subsection{Key Compile-Time Options}

The behavior of KPP is controlled at compile time through \texttt{KPP\_OPTIONS.h}:

\begin{description}
    \item[ALLOW\_KPP] Enable the package (default: ON if pkg/kpp active)
    \item[KPP\_GHAT] Include nonlocal transport coefficient $\gamma$ (default: ON)
    \item[KPP\_SMOOTH\_SHSQ] Horizontally smooth shear squared (default: ON)
    \item[EXCLUDE\_KPP\_SHEAR\_MIX] Turn off shear instability contribution (default: OFF)
    \item[EXCLUDE\_KPP\_DOUBLEDIFF] Turn off double-diffusive enhancement (default: OFF)
    \item[KPP\_SMOOTH\_VISC] Horizontally smooth viscosity field (default: OFF)
    \item[KPP\_SMOOTH\_DIFF] Horizontally smooth diffusivity fields (default: OFF)
\end{description}

\section{Key Parameters and Their Meanings}

\subsection{Boundary Layer Diagnosis}

\begin{description}
    \item[Ricr = 0.3] Critical bulk Richardson number for $h_{bl}$ diagnosis
    \item[minKPPhbl = 1.0] Minimum boundary-layer depth [m]; prevents spurious deepening
    \item[LimitHblStable] If true, limit $h_{bl}$ deepening in stable conditions
\end{description}

\subsection{Interior Mixing Constants}

\begin{description}
    \item[Riinfty = 0.7] Local Richardson number limit for shear instability
    \item[difm0 = 0.005] Background viscosity [m$^2$/s]
    \item[difs0 = 0.005] Background salinity diffusivity [m$^2$/s]
    \item[dift0 = 0.005] Background temperature diffusivity [m$^2$/s]
    \item[difmcon = 0.1] Convective diffusivity (for $N^2 < 0$) [m$^2$/s]
\end{description}

\subsection{Boundary Layer Shape Function}

\begin{description}
    \item[cstar = 10.0] Dimensionless nonlocal transport strength
    \item[conam = 1.8] Momentum boundary layer coefficient
    \item[concs = 1.6] Scalar boundary layer coefficient
\end{description}

\subsection{Turbulent Velocity Scale Lookup}

\begin{description}
    \item[zetam, conc1, ...] Lookup table parameters for unstable conditions (precomputed in \texttt{kpp\_init\_fixed})
    \item[zetas, conc2, ...] Lookup table parameters for stable conditions
\end{description}

\subsection{Numerical Bounds}

\begin{description}
    \item[zmin, zmax] Depth bounds for valid lookup table use [m]
    \item[umin, umax] Velocity bounds [m/s]
    \item[num\_v\_smooth\_Ri] Number of smoothing passes for Richardson number
\end{description}

\section{Initialization Phase}

\subsection{Step 1: kpp\_readparms}

Reads namelist \texttt{KPP\_PARM01} from \texttt{data.kpp}. Sets all runtime parameters and default values. Key actions:

\begin{enumerate}
    \item Open \texttt{data.kpp}
    \item Set defaults for all parameters listed in Section~\ref{sec:kpp_key_params}
    \item Read user-supplied namelist values
    \item Validate parameter consistency
    \item Print configuration summary
\end{enumerate}

Output: All parameters written to common blocks in memory for later access.

\subsection{Step 2: kpp\_init\_fixed}

Computes constants and lookup tables (one-time, called at initialization). Key actions:

\begin{enumerate}
    \item Build turbulent velocity scale lookup tables (\texttt{wmt}, \texttt{wst}) for unstable/stable conditions
    \item Compute derived constants from input parameters
    \item Store vertical grid information from MITgcm
    \item Initialize Richardson number smoothing arrays
\end{enumerate}

Output: Lookup tables and derived constants stored in common blocks for rapid access during runtime.

\subsection{Step 3: kpp\_init\_varia}

Initialize time-varying arrays at model start time. Key actions:

\begin{enumerate}
    \item Initialize \texttt{KPPhbl} to minimum value everywhere
    \item Initialize \texttt{KPPviscAz}, \texttt{KPPdiffKzT}, \texttt{KPPdiffKzS} to background values
    \item Initialize \texttt{KPPghat} (nonlocal transport) to zero
    \item Apply ocean mask based on bathymetry
\end{enumerate}

Output: All KPP diagnostic arrays populated and ready.

\section{Main Runtime Computation: kpp\_calc}

The central routine is \texttt{kpp\_calc(bi, bj, myTime, myIter, myThid)}, called per tile per time step (or at \texttt{kpp\_freq} if not every step).

\subsection{Input Arguments and Implicit Inputs}

\subsubsection{Explicit Arguments}

\begin{itemize}
    \item \texttt{bi, bj} — Tile indices
    \item \texttt{myTime} — Current simulation time (seconds)
    \item \texttt{myIter} — Current iteration number
    \item \texttt{myThid} — Thread ID for parallel execution
\end{itemize}

\subsubsection{Implicit Inputs (via common blocks and module headers)}

\begin{description}
    \item[Ocean State]
    \begin{itemize}
        \item \texttt{THETA}$(i,j,k)$ — Potential temperature [°C]
        \item \texttt{SALT}$(i,j,k)$ — Salinity [psu or g/kg]
        \item \texttt{UVEL}, \texttt{VVEL}$(i,j,k)$ — Velocity components [m/s]
    \end{itemize}
    
    \item[Grid and Geometry]
    \begin{itemize}
        \item \texttt{rC}$(k)$ — Depth of cell centers (negative-downward) [m]
        \item \texttt{drF}$(k)$ — Cell thickness [m]
        \item \texttt{maskC}$(i,j,k)$ — 1 if wet, 0 if land
        \item \texttt{kLowC}$(i,j)$ — Deepest wet level index per column
    \end{itemize}
    
    \item[Surface Forcing]
    \begin{itemize}
        \item \texttt{surfaceForcingU}, \texttt{surfaceForcingV} — Wind stress components [m/s$^2$]
        \item \texttt{surfaceForcingT} — Temperature forcing [K/s]
        \item \texttt{surfaceForcingS} — Salinity forcing [psu/s]
        \item \texttt{Qsw} — Shortwave radiation [W/m$^2$]
    \end{itemize}
    
    \item[Optional Packages]
    \begin{itemize}
        \item Salt plume fields if \texttt{ALLOW\_SALT\_PLUME} compiled
        \item Shortwave 3D fraction if \texttt{SHORTWAVE\_HEATING} and \texttt{KPPuseSWfrac3D}
    \end{itemize}
\end{description}

\subsection{Detailed Computational Sequence}

\subsubsection{Sequence 1: State Conversion (STATEKPP)}

Convert temperature and salinity to KPP working variables:

\begin{enumerate}
    \item Compute in-situ density via equation of state
    \item Compute thermal expansion coefficient $\alpha = -\frac{1}{\rho}\frac{\partial \rho}{\partial T}$
    \item Compute haline contraction coefficient $\beta = \frac{1}{\rho}\frac{\partial \rho}{\partial S}$
    \item Compute local buoyancy difference from surface:
    \begin{equation}
    \Delta b(k) = g[\alpha \Delta T(k) - \beta \Delta S(k)]
    \end{equation}
    \item Compute buoyancy gradient at interfaces (discrete):
    \begin{equation}
    \text{dbloc}(k) = \text{buoyancy difference across level } k
    \end{equation}
\end{enumerate}

Output: Arrays with thermodynamic state converted to buoyancy variables.

\subsubsection{Sequence 2: Surface Forcing (KPP\_FORCING\_SURF)}

Translate MITgcm forcing fields to KPP turbulent scales:

\begin{enumerate}
    \item Compute friction velocity from wind stress:
    \begin{equation}
    u_* = \sqrt[4]{|\boldsymbol{\tau}|^2/\rho_0}
    \end{equation}
    
    \item Compute buoyancy forcing from heat and salt fluxes:
    \begin{equation}
    B_0 = -g(\alpha Q_T + \beta Q_S)/\rho_0
    \end{equation}
    where $Q_T$, $Q_S$ are surface heat and salt fluxes.
    
    \item Compute shortwave-driven buoyancy forcing (if enabled):
    \begin{equation}
    B_{sol} = -g \alpha Q_{sw}/(\rho_0 c_p)
    \end{equation}
    
    \item Compute shear profile relative to surface reference velocity:
    \begin{equation}
    \text{dVsq}(k) = [u_{ref}-u_k]^2 + [v_{ref}-v_k]^2
    \end{equation}
    with reference velocity typically taken from top few levels.
\end{enumerate}

Output: Forcing variables \texttt{ustar}, \texttt{bo}, \texttt{bosol}, \texttt{dVsq} in physical units.

\subsubsection{Sequence 3: Shear Computation}

Compute vertical velocity shear squared:

\begin{equation}
\text{shsq}(k) = \left[\frac{u_k - u_{k+1}}{\Delta z_{k+1/2}}\right]^2 + \left[\frac{v_k - v_{k+1}}{\Delta z_{k+1/2}}\right]^2
\end{equation}

Average from staggered U, V grid to tracer C-grid. Optional horizontal smoothing if \texttt{KPP\_SMOOTH\_SHSQ} enabled.

\subsubsection{Sequence 4: Interior Mixing Computation (Ri\_iwmix)}

Compute interior (non-boundary-layer) mixing for all levels:

\begin{enumerate}
    \item Compute local Richardson number:
    \begin{equation}
    \text{Ri}(k) = \frac{\text{dbloc}(k) \Delta z}{|\text{shsq}(k)| + \phi_\epsilon}
    \end{equation}
    
    \item Compute shear instability enhancement:
    \begin{equation}
    K_{shear}(k) = \text{difm0} \left[1 - \min\left(\frac{\text{Ri}(k)}{\text{Riinfty}}, 1\right)\right]^3
    \end{equation}
    
    \item Compute static instability diffusivity (if $N^2 < 0$):
    \begin{equation}
    K_{conv}(k) = \text{difmcon}
    \end{equation}
    
    \item Sum contributions:
    \begin{equation}
    K_{interior}(k) = \text{difm0} + K_{shear}(k) + K_{conv}(k)
    \end{equation}
\end{enumerate}

Output: Interior diffusivity profile \texttt{diffus}$(\cdot, \cdot, :, 1)$ for momentum and scalars.

\subsubsection{Sequence 5: Boundary Layer Depth Diagnosis (BLDEPTH)}

Diagnose $h_{bl}$ from bulk Richardson number criterion:

\begin{enumerate}
    \item Compute bulk Richardson numerator profile:
    \begin{equation}
    \text{Ritop}(k) = (z_{surface} - z_k) \cdot \text{dbsfc}(k)
    \end{equation}
    
    \item Compute turbulent velocity contribution:
    \begin{equation}
    V_t^2(z) = -z \cdot w_s(z) \sqrt{|N^2(z)|} \cdot V_{tc}
    \end{equation}
    where $V_{tc}$ is a lookup-derived constant.
    
    \item Compute bulk Richardson number profile:
    \begin{equation}
    \text{Ri}_b(z) = \frac{\text{Ritop}(z)}{\text{dVsq}(z) + V_t^2(z)}
    \end{equation}
    
    \item Find boundary layer base where $\text{Ri}_b(h_{bl}) = \text{Ricr}$ via linear interpolation
    
    \item Enforce minimum depth:
    \begin{equation}
    h_{bl} = \max(h_{bl}, \text{minKPPhbl})
    \end{equation}
    
    \item Determine forcing regime:
    \begin{itemize}
        \item Unstable: $B_0 < 0$ (convection/heating)
        \item Stable: $B_0 > 0$ (cooling/evaporation)
    \end{itemize}
    
    \item Determine grid location case (how $h_{bl}$ aligns with grid points)
\end{enumerate}

Output: \texttt{KPPhbl}, \texttt{stable}, \texttt{casea}, \texttt{kbl}, and Richardson number profile.

\subsubsection{Sequence 6: Boundary Layer Mixing (BLMIX)}

Compute boundary layer mixing coefficients using shape functions:

\begin{enumerate}
    \item Determine turbulent velocity scales from forcing and $h_{bl}$:
    \begin{itemize}
        \item If unstable: lookup tables for \texttt{wm} (momentum), \texttt{ws} (scalars)
        \item If stable: direct formula from Monin-Obukhov theory
    \end{itemize}
    
    \item Compute shape function coefficients matching interior mixing at $z = -h_{bl}$
    
    \item For each level $k$ within the boundary layer ($z > -h_{bl}$):
    \begin{equation}
    K(k) = h_{bl} \cdot w(k) \cdot G(k; h_{bl})
    \end{equation}
    where $G$ is the cubic polynomial shape function.
    
    \item Compute nonlocal transport coefficient (if \texttt{KPP\_GHAT} enabled):
    \begin{equation}
    \gamma(k) = -\frac{\text{cstar}}{\int_0^1 G(\sigma)d\sigma} G(\sigma(k))
    \end{equation}
\end{enumerate}

Output: Boundary layer diffusivity profile, \texttt{ghat} nonlocal coefficient.

\subsubsection{Sequence 7: Matching (ENHANCE)}

Smooth the transition between interior and boundary layer mixing:

\begin{enumerate}
    \item Evaluate interior and boundary layer coefficients at $z = -h_{bl}$
    \item Adjust boundary layer shape function polynomial to match
    \item Interpolate smoothly across the boundary-layer base
    \item Ensure no numerical discontinuities
\end{enumerate}

\subsubsection{Sequence 8: Merge Profiles and Export}

For all levels $k$:

\begin{equation}
K_{\text{final}}(k) = \begin{cases}
K_{\text{BL}}(k) & \text{if } k < -h_{bl} \\
\max(K_{\text{BL}}(k), K_{\text{interior}}(k)) & \text{at } k = -h_{bl} \\
K_{\text{interior}}(k) & \text{if } k > -h_{bl}
\end{cases}
\end{equation}

Store in package arrays:
\begin{itemize}
    \item \texttt{KPPviscAz} — Momentum viscosity
    \item \texttt{KPPdiffKzT} — Temperature diffusivity
    \item \texttt{KPPdiffKzS} — Salinity diffusivity
    \item \texttt{KPPghat} — Nonlocal transport (set to zero above $h_{bl}$)
    \item \texttt{KPPhbl} — Boundary layer depth
\end{itemize}

\section{Integration with the Ocean Model}

\subsection{Where KPP Sits in Control Flow}

The main call sequence is:

\begin{enumerate}
    \item \texttt{model/src/do\_oceanic\_phys.F} — Decides whether KPP should run
    \item If active: calls \texttt{kpp\_calc(bi,bj,myTime,myIter,myThid)} for each tile
    \item After tile loops: \texttt{kpp\_do\_exch(myThid)} exchanges halo regions
    \item Later: Model momentum and tracer solvers use \texttt{KPPviscAz}, \texttt{KPPdiffKzT}, \texttt{KPPdiffKzS}
    \item Tracer solvers also apply nonlocal flux via \texttt{KPPghat}
    \item At I/O time: \texttt{kpp\_output()} and diagnostics fill registered fields
\end{enumerate}

The critical division: \texttt{kpp\_calc} computes and stores coefficients, but does not directly update tendencies. Those solvers read the stored arrays.

\subsection{Momentum Mixing: kpp\_calc\_visc}

Called from \texttt{model/src/calc\_viscosity.F} to insert KPP viscosity:

\begin{description}
    \item[Input] Tile indices, level $k$, existing momentum viscosity arrays \texttt{KappaRU}, \texttt{KappaRV}, \texttt{KPPviscAz}
    \item[Output] Updated \texttt{KappaRU}, \texttt{KappaRV}
    \item[Operation] For each level $k$:
    \begin{enumerate}
        \item Copy \texttt{KPPviscAz} from tracer points
        \item Average/interpolate to U and V grid points (staggered)
        \item Insert into momentum viscosity arrays
    \end{enumerate}
\end{description}

\subsection{Tracer Mixing: kpp\_calc\_diff}

Called from \texttt{model/src/calc\_3d\_diffusivity.F}:

\begin{description}
    \item[Input] Tile indices, level (or all levels), tracer identity (THETA, SALT, or passive)
    \item[Output] Updated diffusivity array (\texttt{KappaRT}, \texttt{KappaRS}, or \texttt{KappaPTr})
    \item[Operation] Copy appropriate KPP diffusivity (\texttt{KPPdiffKzT}, \texttt{KPPdiffKzS}, etc.)
\end{description}

Note: Passive tracers use \texttt{KPPdiffKzS} (salinity diffusivity) by default.

\subsection{Nonlocal Transport: kpp\_transport\_*}

Separate routines add the nonlocal flux terms:

\begin{description}
    \item[kpp\_transport\_t] Nonlocal heat flux: $w_e c^* \gamma(k) Q_T$
    \item[kpp\_transport\_s] Nonlocal salt flux: $w_e c^* \gamma(k) Q_S$
    \item[kpp\_transport\_ptr] Nonlocal passive tracer flux
\end{description}

These are crucial for getting the full KPP effect on tracer budgets in unstable conditions.

\section{Diagnostics and Output}

\subsection{Registered Diagnostics}

The package registers the following diagnostic fields via \texttt{kpp\_diagnostics\_init.F}:

\begin{itemize}
    \item \texttt{KPPviscA} — Vertical viscosity [m$^2$/s]
    \item \texttt{KPPdiffT} — Temperature diffusivity [m$^2$/s]
    \item \texttt{KPPdiffS} — Salinity diffusivity [m$^2$/s]
    \item \texttt{KPPghat} — Nonlocal transport coefficient [s/m$^2$]
    \item \texttt{KPPhbl} — Boundary layer depth [m]
    \item \texttt{KPPshsq} — Shear squared $S^2$ [s$^{-2}$]
    \item \texttt{KPPbfsfc} — Surface buoyancy forcing [m$^2$/s$^3$]
    \item \texttt{KPPRi} — Bulk Richardson number [—]
    \item \texttt{KPPdbloc} — Local buoyancy gradient [s$^{-2}$]
    \item \texttt{KPPfrac} — Shortwave fraction heating the BL [—]
\end{itemize}

\subsection{Snapshot Output: kpp\_output}

Called from \texttt{model/src/do\_the\_model\_io.F}:

\begin{description}
    \item[Trigger] Every \texttt{kpp\_dumpFreq} seconds if \texttt{KPPwriteState = .TRUE.}
    \item[Format] MDSIO (MITgcm default binary) or MNC (netCDF)
    \item[Fields] 3D fields for viscosity, diffusivity, boundary layer depth, nonlocal coefficient, etc.
\end{description}

Note: Unlike GGL90, KPP does not require pickup (restart) files because it is diagnostic.

\section{Comparison with GGL90 (TKE Closure)}

\subsection{Fundamental Differences}

\begin{table}[h]
\centering
\caption{KPP vs.\ GGL90 Comparison}
\begin{tabular}{lll}
\toprule
\textbf{Aspect} & \textbf{KPP} & \textbf{GGL90} \\
\midrule
\textbf{Approach} & Diagnostic boundary layer & Prognostic TKE closure \\
\textbf{Prognostic Variables} & None & TKE(z) at all levels \\
\textbf{Time Dependence} & Memoryless & Evolves over time \\
\textbf{History Sensitivity} & No (depends only on current state) & Yes (depends on TKE history) \\
\textbf{Mixing Profile Shape} & Polynomial within BL, interior outside & Smooth, depth-dependent everywhere \\
\textbf{Boundary Layer Depth} & Explicitly diagnosed from Ri criterion & Implicit in TKE decay \\
\textbf{Transient Response} & Instantaneous & Gradual TKE spin-up/down \\
\textbf{Interior Mixing} & Background + shear + static instability & Continuous $S^2, N^2$ dependence \\
\textbf{Surface Forcing Coupling} & Via mixed-layer depth diagnosis & Via turbulent production/dissipation \\
\textbf{Nonlocal Transport} & Explicitly included via $\gamma(z)$ & Not included (optional via IDEMIX) \\
\textbf{Computational Cost} & Lower (diagnostic, no PDE solve) & Higher (solves TKE PDE per column) \\
\textbf{Restart Requirements} & None required & Must store TKE state \\
\end{tabular}
\end{table}

\subsection{Physical Behavior Comparison}

\subsubsection{Wind-Driven Mixing (Stable Stratification)}

\begin{description}
    \item[KPP]
    \begin{itemize}
        \item Wind stress diagnoses boundary layer depth via bulk Richardson number
        \item Boundary layer has polynomial mixing profile (cubic decay)
        \item Transition is sharp at boundary layer base
        \item Interior mixing is weak and constant
    \end{itemize}
    
    \item[GGL90]
    \begin{itemize}
        \item Wind stress generates shear production of TKE
        \item TKE diffuses downward, decaying as $\text{TKE}^{3/2}/L$
        \item Mixing profile is smooth, approximately exponential
        \item Penetration depth depends on TKE balance with dissipation
    \end{itemize}
\end{description}

\subsubsection{Convective Mixing (Unstable Stratification)}

\begin{description}
    \item[KPP]
    \begin{itemize}
        \item Negative buoyancy flux directly increases boundary layer depth
        \item Strong nonlocal transport term $\gamma(z)$ acts in unstable layers
        \item Boundary layer can deepen rapidly
        \item Nonlocal fluxes dominate tracer mixing in deep convection
    \end{itemize}
    
    \item[GGL90]
    \begin{itemize}
        \item Negative $N^2$ allows TKE to grow without stratification limit
        \item Buoyancy forcing directly reduces TKE (negative dissipation contribution)
        \item Mixing extends to depth where TKE dissipates
        \item Can produce very deep mixed layers naturally over days
    \end{itemize}
\end{description}

\subsubsection{Double-Diffusive Layers}

\begin{description}
    \item[KPP]
    \begin{itemize}
        \item Optional double-diffusion routines enhance diffusivity in fingering/diffusive convection
        \item Modifies background profiles before boundary layer is diagnosed
        \item Captures salt-finger structures and stratification changes
    \end{itemize}
    
    \item[GGL90]
    \begin{itemize}
        \item Stratification-dependent mixing length naturally reduces mixing in stable, stratified layers
        \item Double diffusion not explicitly parameterized (could be added)
    \end{itemize}
\end{description}

\subsection{Validation and Tuning Considerations}

\subsubsection{Design Intent}

\begin{itemize}
    \item \textbf{KPP:} Designed for subtropical and tropical oceans with shallow, well-defined boundary layers
    \item \textbf{GGL90:} Designed for energy-consistent full-depth mixing, especially important in deep convection
\end{itemize}

\subsubsection{When to Use KPP vs.\ GGL90}

\begin{description}
    \item[Use KPP for:]
    \begin{itemize}
        \item Process studies with rapid initialization needed
        \item Applications requiring fast execution
        \item Mixed-layer budgets in stable stratification
        \item Comparison with KPP-tuned observations
        \item Diagnostic (stateless) requirements
    \end{itemize}
    
    \item[Use GGL90 for:]
    \begin{itemize}
        \item Long integrations where TKE history matters
        \item Global state estimation (ECCOv4 approach)
        \item Coupled Earth system models
        \item Deep convection events
        \item Studies of internal wave interactions
    \end{itemize}
\end{description}

\section{Python 1D Implementation Status}

The KPP scheme has been faithfully ported to Python 1D for research and validation. Core physics is implemented with comprehensive validation through regression testing. See documentation accompanying the 1D mixing model for complete Python implementation details.

\appendix

\section{Parameter Table}

\begin{longtable}{lp{6cm}ll}
\caption{KPP Parameters: Names, Descriptions, Units, and Typical Values} \\
\toprule
\textbf{Parameter Name} & \textbf{Description} & \textbf{Units} & \textbf{Value} \\
\midrule
\endhead
\bottomrule
\endfoot

Ricr & Critical bulk Richardson number (boundary layer diagnosis) & — & 0.3 \\
Riinfty & Local Richardson number limit (shear mixing) & — & 0.7 \\
minKPPhbl & Minimum boundary layer depth & m & 1.0 \\
\midrule
difm0 & Background viscosity and diffusivity & m$^2$/s & $5 \times 10^{-3}$ \\
difs0 & Background salinity diffusivity & m$^2$/s & $5 \times 10^{-3}$ \\
dift0 & Background temperature diffusivity & m$^2$/s & $5 \times 10^{-3}$ \\
difmcon & Convective instability diffusivity (if $N^2 < 0$) & m$^2$/s & 0.1 \\
\midrule
cstar & Nonlocal transport strength coefficient & — & 10.0 \\
conam & Momentum boundary layer coefficient & — & 1.8 \\
concs & Scalar boundary layer coefficient & — & 1.6 \\
\midrule
zetam, zetad & Lookup table parameters (unstable/stable) & — & Computed \\
conc1-3 & Shape function coefficients & — & Computed \\
\midrule
zmin, zmax & Depth bounds for lookup table & m & Computed from grid \\
umin, umax & Velocity bounds for lookup table & m/s & $\sim 0.0, \sim 0.5$ \\
\midrule
num\_v\_smooth\_Ri & Number of smoothing passes for Richardson number & — & 1 \\
von\_Karman & von Kármán constant & — & 0.4 \\

\end{longtable}

\section{References}

\begin{itemize}
    \item Large, W.\ G., J.\ C.\ McWilliams, and S.\ C.\ Doney (1994). Oceanic vertical mixing: a review and an agenda for future research. In \textit{Proceedings of the WOCE Physical Programme Scientific Workshop No. 6}, held in Melbourne, Australia, Feb.\ 9--11, 1994 (pp.\ 123--172). International WOCE Project Office Publication No.\ 14.
    
    \item Mellor, G.\ L., and T.\ Yamada (1982). Development of a turbulence closure model for geophysical fluid problems. \textit{Reviews of Geophysics and Space Physics}, 20(4), 851–875.
    
    \item Deardorff, J.\ W. (1970). Convective velocity and temperature scales for the unstable planetary boundary layer and for Rayleigh convection. \textit{Journal of the Atmospheric Sciences}, 27, 1211–1213.
    
    \item Adcroft, A., C.\ Hill, and J.\ Marshall (1997). Representation of topography by shaved cells in a height coordinate ocean model. \textit{Monthly Weather Review}, 125, 2293–2315.
    
    \item Marshall, J., A.\ Adcroft, C.\ Hill, L.\ Perelman, and C.\ Heisey (1997). A finite-volume, incompressible Navier-Stokes model for studies of the ocean on parallel computers. \textit{Journal of Computational Physics}, 141, 1–37.
\end{itemize}

\end{document}
